\documentclass[runningheads]{llncs}

\usepackage[T1]{fontenc}
\usepackage[spanish]{babel}
\usepackage{graphicx}
\usepackage{booktabs}
\usepackage{amsmath}
\usepackage{url}
\usepackage{multirow}
\usepackage{array}
\usepackage[hidelinks]{hyperref}

\begin{document}

\title{Sistema de Generaci\'on Aumentada por Recuperaci\'on para
Consultas de Turismo Latinoamericano con B\'usqueda H\'ibrida,
Expansi\'on de Consultas y Fallback Web Autom\'atico}

\titlerunning{Sistema RAG para Turismo Latinoamericano}

\author{Ronald Alfonso}
\authorrunning{R. Alfonso}
\institute{Universidad, Sistemas de Recuperaci\'on de Informaci\'on\\
\email{ronald@bool.mx}}

\maketitle

%-----------------------------------------------------------------------
\begin{abstract}
Se presenta un sistema de Generaci\'on Aumentada por Recuperaci\'on
(RAG) para consultas de turismo en Am\'erica Latina en espa\~nol.
Combina b\'usqueda h\'ibrida BM25~+~KNN sobre OpenSearch con embeddings
densos de la API Mistral (1\,024 dimensiones), expansi\'on de consultas
mediante retroalimentaci\'on de pseudo-relevancia y expansor LLM,
reordenamiento con \textit{cross-encoder}, posicionamiento multifactor
con impulso por pa\'is y retroalimentaci\'on expl\'icita del usuario,
y fallback autom\'atico a Wikivoyage. Se eval\'ua con 30 consultas y
juicios de relevancia manuales usando MAP, MRR, Precisi\'on, Recall y
NDCG en $k\in\{5,10,20\}$. El recuperador base logra MAP~=~0{,}550 y
MRR~=~0{,}887; el pipeline completo obtiene MAP~=~0{,}395 y
MRR~=~0{,}777 con cobertura robusta ante consultas fuera del \'indice.
\keywords{Generaci\'on Aumentada por Recuperaci\'on \and B\'usqueda
H\'ibrida \and Expansi\'on de Consultas \and Reordenamiento \and
Recuperaci\'on de Informaci\'on Tur\'istica \and OpenSearch}
\end{abstract}

%-----------------------------------------------------------------------
\section{Introducci\'on}

La r\'apida proliferaci\'on de modelos de lenguaje de gran escala (LLM)
ha reavivado el inter\'es por anclar sus salidas en corpus verificados.
La Generaci\'on Aumentada por Recuperaci\'on (RAG)~\cite{lewis2020rag}
acopla un generador LLM param\'etrico con un recuperador no param\'etrico,
produciendo respuestas factualmente fundamentadas sin perder la fluidez
de la generaci\'on neuronal.

El turismo latinoamericano presenta retos particulares: consultas muy
espec\'ificas (pa\'ises, ciudades, temporadas, actividades), corpus
heterog\'eneo y escasez de benchmarks en espa\~nol.

Las contribuciones de este trabajo son:
\begin{enumerate}
  \item Arquitectura extremo a extremo de un sistema RAG para turismo
  latinoamericano, desplegado como microservicios Docker Compose.
  \item M\'odulo de posicionamiento multifactor: relevancia sem\'antica,
  impulso por pa\'is detectado y retroalimentaci\'on expl\'icita del
  usuario como factores de ranking.
  \item Dos m\'odulos opcionales completos: expansi\'on de consultas
  (LLM + PRF) y evaluaci\'on cuantitativa con m\'etricas IR.
  \item Benchmark de 30 consultas en espa\~nol con juicios de relevancia.
\end{enumerate}

%-----------------------------------------------------------------------
\section{Trabajo Relacionado}
\label{sec:related}

\textbf{Fundamentos de RAG.}
Lewis et al.~\cite{lewis2020rag} introdujeron el paradigma RAG con DPR.
Asai et al.~\cite{asai2023selfrag} propusieron Self-RAG con tokens de
cr\'itica para recuperaci\'on selectiva.

\textbf{Recuperaci\'on h\'ibrida.}
BM25 disperso combinado con b\'usqueda vectorial densa supera a cada
m\'etodo por separado~\cite{ma2021replication,luan2021sparse}. OpenSearch
soporta pipelines de puntuaci\'on h\'ibrida nativos.

\textbf{Expansi\'on de consultas.}
PRF cl\'asico~\cite{rocchio1971relevance} y enfoques modernos con
LLMs~\cite{gao2022hyde} se combinan en este trabajo.

\textbf{Reordenamiento y RI multiling\"ue.}
\textit{Cross-encoders} MS-MARCO~\cite{reimers2019sentence} y benchmarks
como MIRACL~\cite{zhang2023miracl} son la base de comparaci\'on.

%-----------------------------------------------------------------------
\section{Arquitectura del Sistema}
\label{sec:architecture}

\subsection{Visi\'on General}

Seis contenedores Docker Compose: \textit{crawler} Scrapy,
\textit{procesador}, API FastAPI, frontend Vue~3/Quasar, servidor LLM
Ollama y base de datos vectorial OpenSearch.

\subsection{Ingesta de Datos}

El crawler recopila art\'iculos de turismo latinoamericano. Un script
offline procesa el volcado XML de Wikivoyage en espa\~nol, filtrando
art\'iculos de Am\'erica Latina y el Caribe por coincidencia de pa\'is
en t\'itulo, categor\'ias y texto, con lista de bloqueo para otras regiones.

\subsection{Procesamiento e Indexaci\'on}

El texto se segmenta en \textit{chunks} de~512 tokens con solapamiento
de~50. El \'indice OpenSearch configura el analizador \texttt{spanish}
(tokenizaci\'on, \textit{stopwords} del espa\~nol y \textit{stemming}
Snowball) sobre el campo BM25. En paralelo, cada \textit{chunk} se
representa con \texttt{mistral-embed} (1\,024 dimensiones, L2) para
b\'usqueda KNN con HNSW ($m=16$, $ef\_\mathrm{construction}=128$,
similitud coseno).

\subsection{Pipeline de Recuperaci\'on}

\textbf{Control de dominio.} Clasificador LLM que bloquea el fallback
web para consultas fuera del dominio tur\'istico.

\textbf{Detecci\'on de pa\'is.} Detector basado en reglas para~70 alias
de pa\'ises. El pa\'is detectado se usa como factor de impulso en el
ranking, no como filtro duro.

\textbf{Expansi\'on de consultas.} Expansor LLM (paraf\'raseo) + PRF
(reponderaci\'on de t\'erminos con top-$k$ resultados). Configurables
y combinables.

\textbf{B\'usqueda h\'ibrida.}
$s_{\text{hybrid}} = \alpha\,\hat{s}_{\text{BM25}} +
(1{-}\alpha)\,\hat{s}_{\text{KNN}}$.

\textbf{M\'odulo de posicionamiento multifactor.}
\begin{equation}
  R(q,d) = s_{\text{hybrid}}(q,d)
           \times b_{\text{pa\'is}}(q,d)
           \times b_{\text{feedback}}(d)
  \label{eq:ranking}
\end{equation}
$b_{\text{pa\'is}}=1{,}5$ si el pa\'is de la consulta coincide con el
del documento (1 si no), $b_{\text{feedback}}=1{,}2$ si el documento
tiene retroalimentaci\'on positiva registrada (1 si no calificado).
Con reordenamiento, $s_{\text{hybrid}}$ se sustituye por la puntuaci\'on
del \textit{cross-encoder}
(\texttt{ms-marco-MiniLM-L-6-v2}).

\textbf{Fallback web.} Cascada: \'indice principal $\to$ cach\'e web $\to$
Wikivoyage en vivo. Se activa cuando el \'indice principal devuelve
menos de un umbral m\'inimo de resultados.

\subsection{Retroalimentaci\'on Expl\'icita del Usuario}

La interfaz incluye botones \textit{\'util} / \textit{no \'util} bajo
cada respuesta. Al hacer clic se invoca \texttt{POST /feedback}, que
crea el \'indice \texttt{feedback} en OpenSearch si no existe y
persiste la tupla $(consulta,\allowbreak\;chunk\_id,\allowbreak\;relevante)$.
El factor $b_{\text{feedback}}$ de la Ec.~\ref{eq:ranking} se aplica
en tiempo real a b\'usquedas futuras. El proxy del servidor de
desarrollo Quasar redirige \texttt{/feedback} al backend en el puerto
8000, de forma an\'aloga a las rutas \texttt{/ask} y \texttt{/health}.

\subsection{Generaci\'on de Respuestas}

El top-$k$ de \textit{chunks} con citas se pasa a Ollama
(\texttt{qwen2.5:3b}). La API expone \texttt{/ask} s\'incrono y
\texttt{/ask/stream} SSE.

%-----------------------------------------------------------------------
\section{Configuraci\'on Experimental}
\label{sec:evaluation}

\subsection{Colecci\'on de Prueba}

30 consultas en espa\~nol sobre turismo latinoamericano (playas,
gastronom\'ia, senderismo, sitios culturales, seguridad, transporte)
en~15 pa\'ises. Juicios de relevancia binarios por pooling top-20.

\subsection{Configuraciones}

\textbf{Baseline} (h\'ibrido puro), \textbf{Expansion} (+LLM+PRF),
\textbf{Reranking} (+\textit{cross-encoder}),
\textbf{Expansion+Reranking} y \textbf{Full Pipeline}.

\subsection{M\'etricas}

MAP, MRR, P@$k$, R@$k$, NDCG@$k$ para $k\in\{5,10,20\}$.

%-----------------------------------------------------------------------
\section{Resultados y Discusi\'on}
\label{sec:results}

\begin{table}[t]
\centering
\caption{MAP y MRR (30 consultas). Negrita: mejor valor.}
\label{tab:main}
\begin{tabular}{lcc}
\toprule
\textbf{Configuraci\'on} & \textbf{MAP} & \textbf{MRR} \\
\midrule
Baseline              & \textbf{0{,}550} & \textbf{0{,}887} \\
Expansion             & 0{,}393          & 0{,}737          \\
Reranking             & 0{,}384          & 0{,}667          \\
Expansion+Reranking   & 0{,}350          & 0{,}739          \\
Full Pipeline         & 0{,}395          & 0{,}777          \\
\bottomrule
\end{tabular}
\end{table}

\begin{table}[t]
\centering
\caption{P@$k$ y R@$k$. Negrita: mejor por columna.}
\label{tab:prk}
\setlength{\tabcolsep}{5pt}
\begin{tabular}{lcccccc}
\toprule
\multirow{2}{*}{\textbf{Configuraci\'on}}
  & \multicolumn{3}{c}{\textbf{P@k}} & \multicolumn{3}{c}{\textbf{R@k}} \\
\cmidrule(lr){2-4}\cmidrule(lr){5-7}
  & @5 & @10 & @20 & @5 & @10 & @20 \\
\midrule
Baseline            & \textbf{0{,}620} & \textbf{0{,}487} & \textbf{0{,}307} & \textbf{0{,}380} & \textbf{0{,}601} & \textbf{0{,}750} \\
Expansion           & 0{,}533 & 0{,}400 & 0{,}253 & 0{,}317 & 0{,}479 & 0{,}597 \\
Reranking           & 0{,}400 & 0{,}373 & 0{,}300 & 0{,}208 & 0{,}401 & 0{,}747 \\
Expansion+Reranking & 0{,}440 & 0{,}387 & 0{,}253 & 0{,}243 & 0{,}431 & 0{,}584 \\
Full Pipeline       & 0{,}547 & 0{,}403 & 0{,}253 & 0{,}326 & 0{,}473 & 0{,}582 \\
\bottomrule
\end{tabular}
\end{table}

\begin{table}[t]
\centering
\caption{NDCG@$k$. Negrita: mejor valor.}
\label{tab:ndcg}
\begin{tabular}{lccc}
\toprule
\textbf{Configuraci\'on} & \textbf{@5} & \textbf{@10} & \textbf{@20} \\
\midrule
Baseline            & \textbf{0{,}688} & \textbf{0{,}664} & \textbf{0{,}704} \\
Expansion           & 0{,}563 & 0{,}519 & 0{,}540 \\
Reranking           & 0{,}443 & 0{,}451 & 0{,}580 \\
Expansion+Reranking & 0{,}491 & 0{,}490 & 0{,}519 \\
Full Pipeline       & 0{,}584 & 0{,}511 & 0{,}540 \\
\bottomrule
\end{tabular}
\end{table}

El baseline domina en MAP y MRR. La expansi\'on LLM introduce ruido en
un dominio especializado; el \textit{cross-encoder} ingl\'es degrada el
reordenamiento sobre contenido espa\~nol. El Full Pipeline es
arquitect\'onicamente superior para producci\'on: gestiona consultas sin
cobertura, aplica posicionamiento multifactor y recoge retroalimentaci\'on
del usuario.

%-----------------------------------------------------------------------
\section{Manual de Usuario}
\label{sec:manual}

\subsection{Acceso}

Con el sistema activo (\texttt{docker compose up -d}), acceder a
\url{http://localhost:9000} desde cualquier navegador. No requiere
registro.

\subsection{Descripci\'on de la Interfaz}

La interfaz consta de tres zonas:
\begin{enumerate}
  \item \textbf{\'Area de conversaci\'on}: historial de preguntas y
  respuestas en formato chat. Cada respuesta del asistente muestra el
  texto generado en tiempo real (efecto de escritura token a token) y,
  al finalizar, un chip con el pa\'is detectado si corresponde.
  \item \textbf{Panel de fuentes} (expandible): lista los fragmentos del
  corpus usados, con puntuaci\'on de relevancia, t\'itulo y enlace a la
  p\'agina original. Los fragmentos se ordenan de mayor a menor
  relevancia seg\'un la funci\'on $R(q,d)$ de la Ec.~\ref{eq:ranking}.
  \item \textbf{Barra de entrada}: campo de texto en la parte inferior
  con bot\'on de env\'io.
\end{enumerate}

\subsection{C\'omo Realizar una Consulta}

\begin{enumerate}
  \item Escribir la pregunta en lenguaje natural en espa\~nol, por
  ejemplo: \textit{``¿Qu\'e playas hay en Colombia?''} o
  \textit{``¿Qu\'e ver en Cusco en julio?''}.
  \item Presionar \textit{Enter} o el bot\'on de env\'io.
  \item Esperar la respuesta. El chip de pa\'is confirma si el sistema
  detect\'o un filtro geogr\'afico; el panel de fuentes muestra de d\'onde
  proviene la informaci\'on.
\end{enumerate}

\subsection{Retroalimentaci\'on al Sistema}

Bajo cada respuesta aparecen dos botones: \textbf{\'Util} y
\textbf{No \'util}. Hacer clic en uno env\'ia la valoraci\'on al endpoint
\texttt{POST /feedback}. El bot\'on seleccionado se resalta para confirmar
el registro. Esta retroalimentaci\'on mejora el posicionamiento de
documentos similares en consultas futuras mediante el factor
$b_{\text{feedback}}$.

\subsection{Posicionamiento de Resultados}

Los resultados se ordenan considerando: (1)~relevancia sem\'antica
(embeddings Mistral + BM25 con analizador espa\~nol), (2)~correspondencia
de pa\'is (impulso $\times 1{,}5$ para documentos del pa\'is mencionado),
y (3)~retroalimentaci\'on positiva previa ($\times 1{,}2$). Las fuentes
en el panel expandible reflejan ese mismo orden.

%-----------------------------------------------------------------------
\section{Opini\'on Cr\'itica del Sistema}
\label{sec:critica}

\subsection{Bondades}

\begin{itemize}
  \item \textbf{Arquitectura modular y reproducible.} Cada componente
  corre en un contenedor Docker independiente; el sistema completo se
  levanta con un solo comando y puede reproducirse en cualquier m\'aquina
  con Docker instalado.
  \item \textbf{B\'usqueda h\'ibrida efectiva.} La combinaci\'on de BM25
  con analizador \texttt{spanish} y KNN denso cubre tanto la coincidencia
  l\'exica como la similitud sem\'antica, compensando las limitaciones de
  cada enfoque por separado.
  \item \textbf{Cobertura din\'amica.} El fallback de tres niveles
  garantiza una respuesta aunque el \'indice principal est\'e vac\'io:
  primero busca en cach\'e web y finalmente consulta Wikivoyage en vivo,
  indexando el resultado para reutilizaci\'on futura.
  \item \textbf{Posicionamiento multifactor.} El ranking incorpora
  relevancia sem\'antica, impulso geogr\'afico y retroalimentaci\'on
  expl\'icita del usuario, superando el enfoque de relevancia pura.
  \item \textbf{M\'odulos opcionales funcionales.} Expansi\'on de
  consultas (LLM + PRF) y evaluaci\'on cuantitativa est\'an completamente
  implementados e integrados en el pipeline.
\end{itemize}

\subsection{Insuficiencias y C\'omo Solventarlas}

\begin{itemize}
  \item \textbf{\textit{Cross-encoder} en ingl\'es sobre contenido
  espa\~nol.} El modelo \texttt{ms-marco-MiniLM-L-6-v2} fue entrenado
  sobre consultas inglesas, lo que degrada su puntuaci\'on de relevancia
  sobre texto en espa\~nol. \textit{Soluci\'on}: ajustar un
  \textit{cross-encoder} sobre pares consulta-documento tur\'isticos en
  espa\~nol, o usar \texttt{cross-encoder/mmarco-mMiniLMv2-L12-H384},
  variante multiling\"ue del mismo modelo.
  \item \textbf{Expansi\'on LLM de dominio abierto introduce ruido.}
  El expansor LLM genera par\'afrasis sin restricci\'on de dominio,
  desplazando el centroide sem\'antico de la consulta. \textit{Soluci\'on}:
  incluir un prompt con ejemplos de consultas tur\'isticas para acotar
  el espacio de expansi\'on.
  \item \textbf{Retroalimentaci\'on no actualiza el \'indice vectorial.}
  El factor $b_{\text{feedback}}$ impulsa documentos en tiempo de
  b\'usqueda pero no actualiza los embeddings del corpus. \textit{Soluci\'on}:
  implementar un ciclo de reindexaci\'on peri\'odico que pondere los
  embeddings de documentos con alta valoraci\'on positiva.
  \item \textbf{Popularidad y frescura no consideradas.} El m\'odulo de
  posicionamiento no incorpora popularidad del destino ni frescura del
  documento. \textit{Soluci\'on}: agregar factores $b_{\text{frescura}}$
  basados en \texttt{created\_at} y $b_{\text{popularidad}}$ basado en
  frecuencia de aparici\'on en resultados positivos.
  \item \textbf{Sin m\'odulo multimodal.} El sistema procesa
  \'unicamente texto. \textit{Soluci\'on}: integrar embeddings de
  im\'agenes con CLIP para recuperaci\'on por foto de destino,
  almacen\'andolos como vectores adicionales en el mismo \'indice HNSW.
\end{itemize}

%-----------------------------------------------------------------------
\section{Conclusiones}
\label{sec:conclusions}

Se present\'o un sistema RAG extremo a extremo para turismo latinoamericano
con \'indice h\'ibrido OpenSearch, analizador espa\~nol para BM25,
embeddings Mistral y pipeline modular. El m\'odulo de posicionamiento
combina relevancia, impulso geogr\'afico y retroalimentaci\'on del
usuario. Se implementaron dos m\'odulos opcionales completos: expansi\'on
de consultas y evaluaci\'on cuantitativa.

El baseline h\'ibrido domina en precisi\'on (MAP~=~0{,}550,
P@5~=~0{,}620); el Full Pipeline aporta cobertura robusta
(MAP~=~0{,}395, MRR~=~0{,}777). La principal l\'inea de mejora es ajustar
el \textit{cross-encoder} al dominio tur\'istico en espa\~nol.

%-----------------------------------------------------------------------
\bibliographystyle{splncs04}
\begin{thebibliography}{10}

\bibitem{lewis2020rag}
Lewis, P., et al.: Generaci\'on aumentada por recuperaci\'on para tareas
de PLN. En: NeurIPS, vol.~33, pp.~9459--9474 (2020)

\bibitem{asai2023selfrag}
Asai, A., et al.: Self-RAG: Aprendiendo a recuperar y criticar.
arXiv:2310.11511 (2023)

\bibitem{ma2021replication}
Ma, X., et al.: Estudio de replicaci\'on del DPR. arXiv:2104.05740 (2021)

\bibitem{luan2021sparse}
Luan, Y., et al.: Representaciones dispersas, densas y de atenci\'on.
TACL \textbf{9}, 329--345 (2021)

\bibitem{rocchio1971relevance}
Rocchio, J.J.: Retroalimentaci\'on de relevancia. En: SMART Retrieval
System, pp.~313--323. Prentice-Hall (1971)

\bibitem{gao2022hyde}
Gao, L., et al.: Recuperaci\'on densa de cero disparos.
arXiv:2212.10496 (2022)

\bibitem{naseri2021ceqe}
Naseri, S., et al.: CEQE: Embeddings para expansi\'on de consultas.
En: ECIR 2021, pp.~467--482 (2021)

\bibitem{reimers2019sentence}
Reimers, N., Gurevych, I.: Sentence-BERT. En: EMNLP 2019,
pp.~3982--3992 (2019)

\bibitem{zhang2023miracl}
Zhang, X., et al.: MIRACL: Dataset multiling\"ue. TACL \textbf{11},
1114--1131 (2023)

\bibitem{robertson2009bm25}
Robertson, S., Zaragoza, H.: BM25 y m\'as all\'a. FTIR \textbf{3}(4),
333--389 (2009)

\end{thebibliography}

\end{document}
